%! AP Statistics — Unit 1, Lesson 1.3 — Student Packet Cover Sheet
\documentclass[10pt]{article}
\usepackage{apstats-article}
\usepackage{apstats-boxes}

%=============================================================================
\begin{document}

% ── Full-bleed navy banner ────────────────────────────────────────────────────
\begin{tikzpicture}[remember picture, overlay]
  \fill[navy] (current page.north west)
    rectangle ([yshift=-0.9in]current page.north east);
\end{tikzpicture}
\vspace{-0.8in}

\begin{center}
  {\color{white}\textbf{\LARGE AP Statistics}}\\[4pt]
  {\color{skymid}\large\textbf{Unit 1: Exploring One-Variable Data}}\\[6pt]
  {\color{white}\textbf{\large Lesson 1.3 \quad Representing a Categorical Variable with Tables}}
\end{center}

\vspace{0.22in}

\namedateperiod

\vspace{0.18in}

% ── LEARNING TARGETS ─────────────────────────────────────────────────────────
\begin{learningtargetbox}
\small
\begin{enumerate}[leftmargin=1.5em, itemsep=5pt]
  \item \ldots{} construct a \textit{frequency table} and a
        \textit{relative frequency table} from a list of categorical data.
  \item \ldots{} interpret relative frequencies in context using complete sentences
        that include the variable and the total group size.
  \item \ldots{} explain why relative frequency (not count) must be used to compare
        the distribution of a categorical variable across groups of different sizes.
\end{enumerate}
\end{learningtargetbox}

\vspace{0.14in}

% ── TABLE OF CONTENTS ────────────────────────────────────────────────────────
\begin{tocbox}
\small
\renewcommand{\arraystretch}{1.5}
\begin{tabularx}{\linewidth}{@{} c l X r @{}}
\rowcolor{sky}
\textbf{\#} & \textbf{Component} & \textbf{Description} & \textbf{Score} \\
\midrule
1 & Warm-Up        & Spiral review + tally practice                    & \blank{1.2cm} \\
2 & Guided Notes   & Frequency tables, relative frequency, interpretation & \blank{1.2cm} \\
3 & Group Activity & Build and compare frequency tables                 & \blank{1.2cm} \\
4 & Exit Ticket    & Complete a relative frequency table + write in context & \blank{1.2cm} \\
5 & Homework       & Table construction, interpretation, AP-style practice & \blank{1.2cm} \\
\midrule
  & \multicolumn{2}{r}{\textbf{Total}} & \blank{1.2cm} \\
\end{tabularx}
\end{tocbox}

\vspace{0.14in}

% ── KEEP IN MIND ─────────────────────────────────────────────────────────────
\begin{remindbox}
\small
A frequency table organizes categorical data; a relative frequency table makes it
\emph{comparable}. Counts can be misleading --- a group with 40 walkers out of 50
is very different from 40 walkers out of 400.

\vspace{6pt}
\begin{tabularx}{\linewidth}{@{} >{\centering\arraybackslash\columncolor{goldbg}}p{0.06\linewidth}
                                    X @{}}
\textbf{\large!} &
\textbf{Relative frequencies must sum to 1 (or 100\%).}\enspace
If yours don't, check for a rounding error or a missing category. \\[6pt]
\textbf{\large!} &
\textbf{Always state the context when interpreting.}\enspace
Write \textit{``44\% of the 50 students surveyed prefer to study at home''} ---
not just ``44\%.'' \\[6pt]
\textbf{\large!} &
\textbf{Counts are fine for one group; relative frequency is required
to compare two groups.}\enspace
This principle reappears in Unit 2 with two-way tables. \\[6pt]
\textbf{\large!} &
\textbf{Relative frequency $=$ proportion $=$ probability (eventually).}\enspace
The arithmetic is identical in Unit 4 --- build the habit now. \\
\end{tabularx}
\end{remindbox}

\end{document}
